4. ( 1 pt ) Determine si el siguiente sistema de ecuaciones lineales en \( F=\mathbb{Z}_{2} \), tiene o no solución. De tener solución exhiba el conjunto solución mostrando y justificando cada paso.
\[
\begin{aligned}
x_{1}+x_{2}+x_{3}+x_{4} & =1 \\
x_{1}+x_{2}+x_{3} & =0 \\
x_{3}+x_{4} & =1
\end{aligned}
\] \\

\begin{solution}
Para resolver este sistema de ecuaciones lineales en el campo \( F = \mathbb{Z}_2 \), seguiremos los siguientes pasos:

Escribir el sistema de ecuaciones en forma de matriz aumentada.
Aplicar operaciones elementales de fila para llevar la matriz a su forma escalonada reducida.
Analizar la solución del sistema en base a la matriz en su forma escalonada reducida.

Primero, escribimos el sistema de ecuaciones en forma de matriz aumentada:
\[
\begin{bmatrix}
1 & 1 & 1 & 1 & \vert & 1 \\
1 & 1 & 1 & 0 & \vert & 0 \\
0 & 0 & 1 & 1 & \vert & 1 \\
\end{bmatrix}
\]
Ahora, aplicamos operaciones elementales de fila para encontrar la forma escalonada reducida de la matriz.

Restamos la fila 2 de la fila 1:

\[
\begin{bmatrix}
0 & 0 & 0 & 1 & \vert & 1 \\
1 & 1 & 1 & 0 & \vert & 0 \\
0 & 0 & 1 & 1 & \vert & 1 \\
\end{bmatrix}
\]

Intercambiamos la fila 1 con la fila 3:

\[
\begin{bmatrix}
0 & 0 & 1 & 1 & \vert & 1 \\
1 & 1 & 1 & 0 & \vert & 0 \\
0 & 0 & 0 & 1 & \vert & 1 \\
\end{bmatrix}
\]

Restamos la fila 1 de la fila 2:

\[
\begin{bmatrix}
0 & 0 & 1 & 1 & \vert & 1 \\
1 & 1 & 0 & 1 & \vert & 1 \\
0 & 0 & 0 & 1 & \vert & 1 \\
\end{bmatrix}
\]

Restamos la fila 3 de la fila 2:

\[
\begin{bmatrix}
0 & 0 & 1 & 1 & \vert & 1 \\
1 & 1 & 0 & 0 & \vert & 0 \\
0 & 0 & 0 & 1 & \vert & 1 \\
\end{bmatrix}
\]
Ahora la matriz está en forma escalonada reducida. Esto nos da el siguiente sistema de ecuaciones:
\[
\begin{aligned}
x_3 + x_4 &= 1 \\
x_1 + x_2 &= 0 \\
x_4 &= 1 \\
\end{aligned}
\]
De la tercera ecuación, $x_4=1$.
Sustituyendo $x_4=1$ en la primera ecuación:
\[
x_3 + 1 = 1 \implies x_3 = 0
\]
Sustituyendo $x_1+x_2=0$:
Podemos elegir $x_1=t$ y $x_2=t$, donde $ t \in \mathbb{Z}_2 $, pues cuales sea dos elementos iguales de este campo al sumarse dan cero.
Por lo tanto, el conjunto solución es:
\[
\begin{aligned}
x_1 &= t \\
x_2 &= t \\
x_3 &= 0 \\
x_4 &= 1 \\
\end{aligned}
\]
donde $ t \in \mathbb{Z}_2 $. Este sistema tiene solución y el conjunto solución depende de t, explícitamente las soluciones son.
\[
\left\{\, (0,\ 0,\ 0,\ 1),\ \ (1,\ 1,\ 0,\ 1) \,\right\}
\]
\end{solution}